
% \section{Methodology}\label{sec:Method}
% We formulate the resource allocation problem as a two-stage stochastic programming problem, where the first stage involves decisions on resource allocation and the second stage represents the realization of uncertain parameters~\parencite{smith2005,jones2010,brown2015}. The objective is to minimize the total cost of relief operations, including procurement, transportation, and distribution costs, subject to various constraints, such as capacity constraints, demand satisfaction requirements, and budget constraints.

% The cost function $C(x)$ is defined as the sum of procurement, transportation, and distribution costs:
% \begin{equation}
% C(x) = \sum_{i=1}^{n} \Biggl(c_i \hat{\mathbf{y}}_i + \sum_{j=1}^{m} d_{ij} t_{ij} + \sum_{k=1}^{l} s_k q_k\Biggr)
% \end{equation}

% We model uncertainty using scenario-based stochastic programming, where multiple scenarios representing different realizations of uncertain parameters are considered. We use historical data, expert opinions, and probabilistic forecasts to generate scenario sets that capture the range of possible outcomes. The stochastic programming model then generates optimal resource allocation decisions that minimize the expected total cost across all scenarios while ensuring robustness against uncertainty.

% The objective function of the stochastic programming model is formulated as follows:
% \begin{equation*}
% \min_{x,y} \quad \sum_{s \in S} \left[ f(x, y, s) \cdot \mathbb{P}(s) \right]
% \end{equation*}
% subject to:
% \begin{align}
% g(x, y, s) \leq 0 &\quad \forall s \in S\\
% h(x, y, s) = 0 &\quad \forall s \in S
% \end{align}

% \begin{theorem}[Optimality Conditions]\label{thm:Opt}
% Let $x^*$ be an optimal solution to the stochastic programming problem. If the objective function and constraint functions are convex, then $x^*$ satisfies the Karush-Kuhn-Tucker (KKT) conditions.
% \end{theorem}

% \begin{proof}{Proof}
% The proof follows from the convex optimization theory, which states that for convex objective and constraint functions, the KKT conditions are necessary and sufficient for optimality. Therefore, if $x^*$ is an optimal solution, it must satisfy the KKT conditions.\Halmos
% \end{proof}

% \begin{algorithm}
% \caption{Random Forest Training}
% \vskip6pt
% \begin{algorithmic}
% \Procedure{RandomForest}{$X_{\text{train}}, y_{\text{train}}, \text{num\_trees}$}
%     \State forest $\gets []$
%     \For{$i \gets 1$ \textbf{to} $\text{num\_trees}$}
%         \State $X_{\text{sampled}}, y_{\text{sampled}} \gets \text{bootstrap\_sample}(X_{\text{train}}, y_{\text{train}})$
%         \State tree $\gets$ \Call{DecisionTree}{$X_{\text{sampled}}, y_{\text{sampled}}$}
%         \State \textbf{append} tree to forest
%     \EndFor
%     \State \textbf{return} forest
% \EndProcedure
% \Statex
% \end{algorithmic}
% \end{algorithm}

%  \begin{figure}
%      \FIGURE
%  %    {\includegraphics{figure-filename.pdf}} %Callout External Image
%  %	Include LaTeX figures
%  {\begin{tikzpicture}[
%   grow=down,
%   level 1/.style={sibling distance=2.5cm, level distance=2.5cm},
%   edge from parent/.style={draw,-latex},
%   every node/.style={font=\footnotesize},
%   decision/.style={shape=rectangle,draw,fill=blue!20},
%   leaf/.style={shape=rectangle,draw,fill=green!20}
% ]
% \node [decision] {Sensor A}
%   child { node [leaf] {Normal} }
%   child { node [decision] {Sensor B}
%     child { node [leaf] {Failure} }
%     child { node [leaf] {Normal} }
%   };
% \end{tikzpicture}}
% {Text of the Figure Caption. \label{fig1}}
% {Text of the notes.}
% \end{figure}

% \begin{algorithm}
% \begin{algorithmic}[1]
% \caption{Random Forest Training II}
% \Procedure{DecisionTree}{$X, y$}
%     \If{\textsc{StoppingCondition}$(X, y)$}
%         \State \textbf{return} LeafNode$(y)$
%     \Else
%         \State $(\text{feature}, \text{threshold}) \gets \textsc{FindBestSplit}(X, y)$
%         \State $\textit{left\_indices} \gets X[\text{feature}] \leq \text{threshold}$
%         \State $\textit{right\_indices} \gets X[\text{feature}] > \text{threshold}$
%         \State $\textit{left\_subtree} \gets$ \Call{DecisionTree}{$X[left\_indices], y[left\_indices]$}
%         \State $\textit{right\_subtree} \gets$ \Call{DecisionTree}{$X[right\_indices], y[right\_indices]$}
%         \State \textbf{return} TreeNode$(\text{feature}, \text{threshold}, left\_subtree, right\_subtree)$
%     \EndIf
% \EndProcedure
% \Statex

% \Procedure{Predict}{$\text{forest}, x$}
%     \State predictions $\gets []$
%     \For{$\text{tree}$ \textbf{in} $\text{forest}$}
%         \State \textbf{append} \Call{TreePredict}{$\text{tree}, x$} to predictions
%     \EndFor
%     \State \textbf{return} \Call{AggregatePredictions}{predictions}
% \EndProcedure
% \Statex

% \Procedure{TreePredict}{$\text{tree}, x$}
%     \If{$\text{tree}$ is a leaf node}
%         \State \textbf{return} $\text{tree.value}$
%     \Else
%         \If{$x[\text{tree.feature}] \leq \text{tree.threshold}$}
%             \State \textbf{return} \Call{TreePredict}{$\text{tree.left\_subtree}, x$}
%         \Else
%             \State \textbf{return} \Call{TreePredict}{$\text{tree.right\_subtree}, x$}
%         \EndIf
%     \EndIf
% \EndProcedure
% \end{algorithmic}
% \end{algorithm}

% \begin{lemma}[Feasibility of Scenario Sets]\label{lem:FSS}
% Given a set of scenarios $S$ generated from historical data and probabilistic forecasts, the scenario set $S$ is feasible if it captures the range of possible outcomes with sufficient coverage.
% \end{lemma}

% \begin{proof}{Proof}
% The proof follows from the definition of feasibility, which requires that the scenario set $S$ includes a representative sample of possible outcomes. Feasibility ensures that the stochastic programming model adequately represents the uncertainty in the problem domain and provides meaningful solutions.\Halmos
% \end{proof}

% \begin{remark}
% This stochastic programming model assumes that demand and supply parameters follow known probability distributions.
% \end{remark}

% \begin{definition}
% A feasible solution to the resource allocation problem satisfies all constraints and requirements without violating any constraints.
% \end{definition}



\section{Detection module}
\label{sec:Method}

\subsection{Graph Construction}
% Why we represent the transaction records as a set of sub-graph?
Financial fraud, particularly money laundering~\parencite{chen2018machine}, poses a major risk to the stability and integrity of financial systems. A fundamental challenge in detecting money laundering is that transactional data are fragmented across banks due to privacy concerns, preventing the construction of a global view of illicit activities~\parencite{le2018preventing, van2024privacy}. This bank isolation limits existing detection systems, leading to both false positives and false negatives, as suspicious behaviors often only emerge when transactions across multiple institutions are jointly analyzed~\parencite{wang2024temporal, amoako2025exploring}. However, directly aggregating transaction data across banks is infeasible due to privacy and regulatory constraints~\parencite{thommandru2023exploring}. 

Traditional machine learning approaches typically treat each transaction as an independent observation and largely ignore the relationships among consecutive money-laundering transactions made by the same group of accounts~\parencite{wang2024temporal, amoako2025exploring, pettersson2022combating, alotibi2022money, jensen2023fighting, lokanan2024predicting}. Such independence assumptions limit their ability to recognize coordinated or evolving laundering schemes, where illicit funds are often transferred through a sequence of interconnected entities to obscure their origin and ownership. Inspires by this observation, we represent the transaction records as a set of graphs in which nodes represent accounts and edges denote fund transfers between them. The conversion to graph-structured data allows us to capture the interdependencies among transactions both direct and indirect within and across laundering groups. 


% Consider a transaction network\footnote{we use graph and network exchangeablely.} $\mathcal{G}=\{\mathcal{G}_1, \mathcal{G}_2, ..., \mathcal{G}_m\}$ consisting of $m$ different subgraphs, where each subgraph $\mathcal{G}_i=\{\mathcal{V}_i, \mathcal{E}_i, X_i\}$ represents a subset of transaction network belonging to a country. The set $\mathcal{V}_i$ denotes the entities (accounts) associated with bank $i$, and $\mathcal{E}_i \subseteq \mathcal{V}_i \times \mathcal{V}_i$ represents the transaction links among these entities. Each entity is further characterized by an attribute matrix $X_i$, which captures its observable financial or behavioral features. Given the adjacency matrix $A_i$,  a graph neural network (GNN)–based encoder $f_\theta$~\parencite{kipf2016semi, DBLP:conf/iclr/XuHLJ19} maps the pair $(A_i, X_i)$ into a set of latent representations $H_i=f_\theta(A_i, X_i)$. These latent embeddings encapsulate the structural and transactional characteristics of each bank’s network and provide the analytical foundation for identifying suspicious interaction patterns and guiding compliance decisions. Our goal is to identify the money-laundering behaviors (i.e., malicious transaction links $\mathcal{E}_i$ among the accounts) across different countries.
% Consider a transaction network\footnote{We use the terms \textit{graph} and \textit{network} interchangeably.} $\mathcal{G}=\{\mathcal{G}_1, \mathcal{G}_2, \ldots, \mathcal{G}_m\}$ consisting of $m$ subgraphs, where each subgraph $\mathcal{G}_i=\{\mathcal{V}_i, \mathcal{E}_i, X^V_i, X^E_i\}$ represents a portion of the global transaction system corresponding to a specific country. The set of nodes $\mathcal{V}_i$ denotes the entities (accounts) within country $i$, and $\mathcal{E}_i \subseteq \mathcal{V}_i \times \mathcal{V}_j$ represents the edge set of transaction links among these entities. Notice that for an edge $\mathcal{E}_i \subseteq \mathcal{V}_i \times \mathcal{V}_j (i\neq j)$, it suggests that this is a cross-border transaction. Each entity is further described by an node-attribute matrix $X^V_i$, which captures its observable financial or behavioral characteristics and each edge is described by an edge-attribute matrix $X^E_i$, which depicts the features of an transaction, such as payment amount, payment method.
Consider a transaction network\footnote{We use the terms \textit{graph} and \textit{network} interchangeably.} $\mathcal{G}=\{\mathcal{G}_1, \mathcal{G}_2, \ldots, \mathcal{G}_m\}$ consisting of $m$ subgraphs, where each subgraph $\mathcal{G}_i=\{\mathcal{V}_i, \mathcal{E}_i, X_i^V, X_i^E, Y_i\}$ represents the financial transaction system within a specific country. The set of nodes $\mathcal{V}_i (|\mathcal{V}_j|=n_i)$ consists of $n_i$ entities (such as individual or corporate accounts) operating in country $i$, while the set of edges $\mathcal{E}_i \subseteq \mathcal{V}_i \times \mathcal{V}_j $ captures the transaction links among these entities. When an edge connects $\mathcal{V}_i$ and $\mathcal{V}_j$ with $i \neq j$, it represents a cross-border transaction, reflecting international fund transfers between entities in different jurisdictions. $Y_i$ is the label vector, where $Y_{i,j}=1$ indicates that the $j$-th transaction is a money-laudering behavior and $Y_{i,j}=0$, otherwise. Here, we convert the set of edges $\mathcal{E}_i$ into an adjacency matrix $A_i\in \mathbb{R}^{n_i \times n_i}$, where $A_i(j,k)=1$ indicates an existing transaction link and  $A_i(j,k)=0$, otherwise. Each node is associated with an attribute matrix $X_i^V\in \mathbb{R}^{n_i \times d_i}$ that records its observable financial or behavioral characteristics, such as transaction frequency, while each edge is characterized by an attribute matrix $X_i^E$ that describes the properties of a transaction, including payment amount, method, and timing. Ultimately, our objective is to identify money-laundering activities, represented by malicious transaction links $\mathcal{E}_i$ among accounts, and to understand how these illicit behaviors manifest and propagate across different countries.

\subsection{Federated Graph-based Detection module}
\label{sec:detection_model}
Given the adjacency matrix $A_i$, we employ a graph neural network (GNN)–based encoder~\parencite{kipf2016semi, DBLP:conf/iclr/XuHLJ19} to map the pair $(A_i, X_i^V)$ into a set of latent representations at the $L$-layer:
\begin{equation}
    H_i^{L+1} = \sigma(A_i H_i^L W^L),
\end{equation}
where $H_i^0=X_i^V$ is the raw node attribute at the first layer, $W^L$ is the weight matrix at the $L$-th layer, and $\sigma(a)=\max(0,a)$ is the ReLU activation function. Intuitively, this graph encoder aggregates information from each entity’s local neighborhood, allowing the representation of an account to be informed not only by its own attributes but also by the behaviors and characteristics of the entities it transacts with. This relational learning process captures how risk exposure or suspicious activity propagates through the transaction network. In financial systems, money-laundering behaviors often span multiple accounts and evolve through multi-step transaction chains. To emulate such interlinked and layered behaviors, we concatenate the latent representations learned from all layers to form the final embedding $H_i = [H_i^0, H_i^1, \ldots, H_i^L]$. This multi-level aggregation allows the module to integrate both local and global structural information, thereby capturing higher-order dependencies that reflect complex transaction patterns, such as circular money flows, bursty transfers, or structured layering activities commonly used to disguise illicit origins. 
% The resulting latent embeddings $H_i$ provide a compact yet comprehensive characterization of both the structural and transactional dynamics of each country’s financial network. 
To transform node-level representations into edge-level embeddings for transaction classification, we design an edge construction function that aggregates the features of the source and destination nodes together with the edge-specific attributes before computing the prediction via a multi-layer perceptron (MLP). Specifically, given two latent representations $h_i \in \mathbb{R}^d$, $h_j \in \mathbb{R}^d$ for nodes $i$, $j$, and the edge $\mathcal{E}_{i,j}$ associated attribute vector $X^E_{ij} \in \mathbb{R}^k$, we construct the edge representation $z_{ij}$ as
\begin{equation}
    z_{ij} = \left[\, h_i \,\|\, (h_j - h_i) \,\|\, X^E_{ij} \,\right],
\end{equation}
where $\|$ denotes the vector concatenation operation. In the representation $z_{ij}$, $h_i$ encodes the intrinsic behavioral and structural characteristics of the source account; $(h_j - h_i)$ models the directional discrepancy between the destination and source accounts, thereby encoding the asymmetric flow of information or funds, which is crucial for detecting suspicious transfers; $X^E_{ij}$ preserve contextual information specific to the transaction, such as the amount or transaction type, providing direct evidence for abnormal activity. Overall, the resulting vector $z_{ij}$ integrates both node-level and edge-level information into a unified representation that characterizes each transaction from multiple perspectives. This representation is then passed through a multi-layer perceptron $f_{\theta}(\cdot)$ followed by a sigmoid activation to predict the probability of the transaction being illicit:
\begin{equation}
    \hat{y}_{ij} = \sigma ( f_\theta(z_{ij})),
\end{equation}
where $\sigma(\cdot)$ denotes the sigmoid function. The model can be trained end-to-end using the binary cross-entropy loss between the predicted probability $\hat{y}_{ij}$ and the ground-truth label $y_{ij}$:
\begin{equation}
    \mathcal{L}_{\text{BCE}} = - \frac{1}{|\mathcal{E}|} \sum_{(i,j) \in \mathcal{E}} \Big[\, y_{ij} \log(\hat{y}_{ij}) + (1 - y_{ij}) \log(1 - \hat{y}_{ij}) \,\Big].
\end{equation}
Through this formulation, the model learns to distinguish illicit transactions by jointly considering node characteristics, their directional interactions, and edge-level contextual evidence. 
However, in real-world financial networks, the distribution of transactions is highly skewed as the vast majority are legitimate, and only a minute fraction correspond to money-laundering activities. This extreme class imbalance poses a fundamental learning challenge because the model tends to be biased toward predicting normal behavior, as minimizing cross-entropy naturally favors the dominant class. Consequently, even a trivial classifier that labels every transaction as normal can achieve deceptively high accuracy but fail to identify the rare yet critical illicit links that drive regulatory risk. To mitigate this imbalance, we adopt the focal loss~\parencite{lin2017focal}, which dynamically reweights training samples to focus more on hard-to-classify, minority cases. The focal loss modifies the cross-entropy objective as follows:
\begin{align}
\label{focal_loss}
    \mathcal{L}_i^{FL} = - \sum_j \big[\alpha(1 - \hat{\mathbf{y}}_j)^{\gamma} y_j \log(\hat{\mathbf{y}}_j) + (1-\alpha)\hat{\mathbf{y}}_j^{\gamma} (1-y_j) \log(1-\hat{\mathbf{y}}_j)\big],
\end{align}
where $\alpha \in [0,1]$ is a class-balancing factor that adjusts the relative importance of the minority (suspicious) versus majority (normal) classes, and $\gamma > 0$ is a focusing parameter that controls the attention on misclassified or uncertain samples. Intuitively, $\alpha$ serves as an ``attention weight'' for class imbalance. A higher $\alpha$ increases the importance of correctly identifying suspicious transactions, ensuring that rare but high-impact money-laundering cases receive more focus, whereas a lower $\alpha$ reduces sensitivity to the minority class. The focusing parameter $\gamma$ emphasizes hard-to-classify samples. For a positive example ($y_j = 1$), when the model is confident (\(\hat{\mathbf{y}}_j \approx 1\)), the factor $(1-\hat{\mathbf{y}}_j)^\gamma$ downweights the loss contribution, reducing the influence of easy cases. Conversely, for uncertain or misclassified positive examples (\(\hat{\mathbf{y}}_j \ll 1\)), the loss remains high, drawing attention to these high-risk transactions. Similarly, for negative examples ($y_j = 0$), the factor $\hat{\mathbf{y}}_j^\gamma$ ensures the model focuses on ambiguous normal transactions. Collectively, $\alpha$ and $\gamma$ enable the model to prioritize learning from rare and high-risk transactions without being overwhelmed by the abundant routine activity.

From a managerial perspective, these hyperparameters encode distinct dimensions of institutional risk tolerance. The class-balancing factor $\alpha$ reflects the strategic trade-off between compliance and customer experience: higher $\alpha$ prioritizes catching illicit transactions even at the cost of more false positives, while lower $\alpha$ minimizes unnecessary account friction to preserve customer relationships. The focusing parameter $\gamma$ controls operational attention allocation: higher $\gamma$ concentrates learning on ambiguous, hard-to-classify transactions—precisely the borderline cases that consume investigator time—while lower $\gamma$ treats all misclassifications more uniformly. This parameterization allows institutions to calibrate detection systems to their specific regulatory environment and operational capacity.

As discussed earlier, one of the most critical challenges in cross-border money-laundering detection lies in the bank-isolation problem, where financial institutions are prohibited from directly sharing transactional data across borders due to strict regulatory and privacy constraints. This limitation hinders the ability to identify complex laundering schemes that intentionally distribute suspicious transfers across multiple banks or jurisdictions. To address this challenge while maintaining data confidentiality, we introduce the concept of a \textit{virtual super-node} $\mathcal{V}^S_{uv}$ ($u \in \mathcal{V}_i, v \in \mathcal{V}_j$), which abstractly represents the transactional relationship between two accounts belonging to different financial institutions. The super-node serves as a ``boundary connector'' between two subgraphs $\mathcal{G}_i$ and $\mathcal{G}_j$, enabling institutions to share structured insights rather than raw transaction details. Formally, we denote by $\mathcal{V}^S$ the set of virtual super-nodes, i.e., $\mathcal{V}^S_{uv} \in \mathcal{V}^S$. Each super-node summarizes localized transaction behavior in a privacy-preserving form that can be exchanged for collaborative analysis. 
% \lc{mention high-order statistical moment instead of just showing the first-order moment (mean aggregation).} 
Specifically, we define the representation of a super-node as: 
\begin{equation}
    h_{\mathcal{V}^S_{iu}} = \frac{1}{|N_{iu}|}\sum_{v \in N_{iu}} h_v, \quad
    h_{\mathcal{V}^S_{jv}} = \frac{1}{|N_{jv}|}\sum_{u \in N_{jv}} h_u,
\end{equation}
where $h_{\mathcal{V}^S_{iu}}$ and $h_{\mathcal{V}^S_{jv}}$ capture aggregated local neighbor information from their neighbor set $N_{iu}$ and $N_{jv}$ in subgraphs $\mathcal{G}_i$ and $\mathcal{G}_j$, respectively. Notice that we can extend the mean aggregation (first-order statistical moment) to the combination of first-order and high-order moment, such as the aggregation of mean and variance, while our empirical results in Section~\ref{supernode_high_order_results} only show the marginal effect of including high-order moment. These embeddings act as compact statistical summaries of local cross-border activities, allowing each institution to participate in a federated learning process without revealing any sensitive data. To ensure semantic consistency across institutions, we introduce a self-consistency loss:
\begin{equation}
    \mathcal{L}_i^{SC} = \sum_{\mathcal{V}^S_{uv} \in \mathcal{V}^S_i} \| h_{\mathcal{V}^S_{iu}} - h_{\mathcal{V}^S_{jv}} \|_2^2,
\end{equation}
where $\mathcal{V}^S_i$ is the subset of super-nodes associated with $\mathcal{G}_i$. This objective encourages alignment between the embeddings of shared super-nodes across banks, ensuring that the same cross-institutional transaction pattern is interpreted consistently. Conceptually, it enables a federated network of banks to ``agree'' on the structural signature of suspicious behavior, thereby strengthening the global defense against coordinated money-laundering schemes while respecting jurisdictional boundaries.


Beyond individual transaction links, money-laundering often manifests in higher-order structures, such as rings, chains, or star-shaped hubs that circulate illicit funds across multiple accounts. To explicitly encourage the model to capture such patterns, we employ a soft-membership differentiable graph cut loss~\parencite{nazi2019gap}, which promotes the identification of dense, irregular clusters in a fully differentiable manner:
\begin{equation}
\mathcal{L}_i^{GCL} = -\frac{\mathrm{trace}(\mathcal{M}_i^\top A_i \mathcal{M}_i)}{\mathrm{trace}(\mathcal{M}_i^\top D_i \mathcal{M}_i)} + \beta \|\mathcal{M}_i^\top \mathcal{M}_i - I\|_F^2,
\end{equation}
where $\mathcal{M}_i = f_\gamma(A_i, H_i)$ represents the soft community assignment of nodes in subgraph $\mathcal{G}_i$, $D_i$ is the diagonal degree matrix of $A_i$, $I$ is the identity matrix, and $\beta$ is a regularization hyperparameter. The first term encourages nodes with strong transactional connections to be grouped together in the same cluster, effectively capturing tightly knit substructures that may indicate coordinated money-laundering activity. The second term regularizes the assignments to promote orthogonality among different clusters, ensuring that each node’s membership is distributed consistently and that communities remain distinct. Intuitively, this loss guides the embeddings to highlight anomalous clusters, such as small rings or star-shaped hubs, that deviate from normal transaction patterns, surfacing laundering groups that may evade detection through standard anomaly scoring alone. Finally, each bank optimizes a multi-objective loss: 
\begin{align}
\label{overall_loss}
    \mathcal{L}_i &= \mathcal{L}_i^{FL} + \lambda_1 \mathcal{L}_i^{GCL} + \lambda_2 \mathcal{L}_i^{SC},
\end{align}
where $\mathcal{L}_i^{FL}$ is the focal loss for edge prediction under weak supervision to distinguish suspicious transaction links, $\mathcal{L}_i^{GCL}$ is the soft-membership differentiable graph cut loss that encourages discovery of dense and anomalous clusters indicative of coordinated laundering schemes, and $\mathcal{L}_i^{SC}$ is the cross-bank consistency loss that aligns shared super-node representations. The hyperparameters $\lambda_1$ and $\lambda_2$ balance the relative importance of the community structure and cross-bank alignment objectives, allowing banks to jointly learn risk-sensitive embeddings while preserving privacy and regulatory compliance.

% Building upon the federated learning paradigm, each bank minimizes its local objective to update parameters $\theta_i$, and periodically transmits only the model updates and super-node embeddings to a central server. Following the federated averaging protocol~\parencite{zhang2021survey, li2020review}, the global parameters are updated as
% \begin{equation}
%     \theta \leftarrow \frac{1}{M} \sum_{i=1}^{M}\theta_i
% \end{equation}
% The updated global parameters $\theta$ are redistributed to all banks, enabling collaborative fraud detection without raw data sharing. 

Building upon the federated learning paradigm, each bank optimizes its local objective to update its model parameters $\theta_i^{local}$ using only its internal transaction data. Critically, raw data never leaves the institution, preserving privacy and regulatory compliance. At predefined communication rounds, each bank transmits only its model updates (gradients or parameters) and the aggregated super-node embeddings to a central server, which acts as a coordination hub rather than a data repository. Following the widely used federated averaging protocol~\parencite{zhang2021survey, li2020review}, the server updates the global model parameters as
\begin{equation}
    \theta^{global} \leftarrow \frac{1}{m} \sum_{i=1}^{m} \theta_i^{local},
\end{equation}
where $m$ is the total number of participating banks. This weighted averaging effectively combines the knowledge learned from heterogeneous local datasets while mitigating biases that may arise from imbalanced or institution-specific transaction patterns. The updated global parameters $\theta^{global}$ are then redistributed to all participating banks, allowing each institution to benefit from the collective intelligence of the network without ever accessing another bank’s raw transactions. Intuitively, this collaborative process enables banks to detect complex, distributed money-laundering schemes that span multiple institutions or jurisdictions. By leveraging federated learning, the system simultaneously improves detection accuracy, respects data privacy, and aligns with regulatory constraints, creating a practical and scalable framework for cross-border financial crime prevention.



\section{Identifying Group Patterns with Cross-bank Personalized PageRank}
\subsection{Personalized PageRank}
% \lc{To do list: 1. Add the intuition of PPR by mentioning the idea of PPR in webpage search. 2. Label propagation. }
While graph-based detection modules achieve strong predictive performance in identifying suspicious transactions, their complexity often limits interpretability, making it challenging for compliance officers to understand why certain accounts are flagged. In our proposed method, we employ a differentiable graph cut loss to capture higher-order structures, such as rings, chains, or star-shaped hubs that circulate illicit funds across multiple accounts. However, this approach alone may not capture group-level suspicious behaviors spanning multiple countries or accounts that often represent a mixture of incomplete laundering schemes rather than compact local structures. To enhance interpretability at the group-level, we incorporate Personalized PageRank (PPR)~\parencite{page1999pagerank}, a well-established algorithm originally developed and deployed by Google in the early 2000s to rank and organize web pages in its search engine. In the web context, PPR models how user interest in a particular page propagates through hyperlinks to identify other closely related pages. This intuitive notion of relevance diffusion motivates its application in financial networks. Specifically, starting from known high-risk accounts, PPR enables us to trace and quantify how suspicion propagates through transactional links to other closely connected accounts. By diffusing risk signals along multi-hop transaction paths, PPR reinforces model predictions while producing transparent and interpretable local clustering results that highlight intermediary accounts, relational dependencies, and latent laundering structures. This relational perspective allows compliance officers to better understand the contextual basis of flagged accounts, thereby increasing trust in model outputs and supporting more targeted and actionable investigations.


% which provides a transparent and intuitive mechanism for tracing and quantifying the relational influence of known high-risk accounts throughout the transaction network. By diffusing suspicion scores along transactional paths, PPR reinforces model predictions while generating interpretable local clustering results that highlight multi-hop relationships, intermediary nodes, and hidden laundering structures. This allows compliance officers to understand the relational context of flagged accounts, increasing trust in model outputs and facilitating actionable investigations. 

Formally, let the transactional network in a country be represented as a directed graph $\mathcal{G}_i$, with a weighted adjacency matrix $P$, where $P_{ij}$ reflects the normalized transaction intensity (e.g., frequency, volume, or dollar amount) from node $i$ to node $j$. Traditional rule-based systems often focus on local anomalies, such as unusually large transfers or sudden spikes in activity. In contrast, money-laundering schemes are inherently relational, spanning multiple transactions, intermediaries, and jurisdictions. PPR addresses this limitation by diffusing initial suspicion scores across the network topology, allowing latent patterns, collusive groups, and hidden intermediaries to surface naturally.
% While the graph-based detection models achieve strong predictive performance in detecting suspicious transactions, their complexity often limits interpretability, making it difficult for compliance officers to understand why certain accounts are flagged. In our proposed method, though we involve the differentiable graph cut loss to capture the higher-order structures,
% such as rings, chains, or star-shaped hubs that circulate illicit funds across multiple accounts, it cannot the captured high-order structure, such as a group of global suspicious accounts sharing the similar patterns (which turns out to be a mixture of multiple groups of incomplete money-laundering behaviors) rather than a compact local money-laundering behaviors.
% To fully capture the group-wise interpretable money-laundering behaviors, we incorporate Personalized PageRank (PPR)~\parencite{}, which provides a transparent mechanism to trace and quantify the relational influence of known high-risk accounts throughout the transaction network. By diffusing suspicion scores along transactional paths, PPR not only reinforces model predictions but also generates interpretable explanations that highlight multi-hop relationships, intermediary nodes, and hidden laundering structures, thereby enhancing trust and facilitating actionable investigation.
% Let the transaction network be represented as a directed graph $\mathcal{G}_i$, we construct a weighted adjacency matrix $W$, where $W_{ij}$ reflecting the normalized transaction intensity (e.g., frequency, volume, or dollar amount) from node $i$ to node $j$. Traditional rule-based systems often focus on local anomalies, such as unusually large transfers or sudden account activity. However, money-laundering schemes are inherently relational, spanning multiple transactions, intermediaries, and jurisdictions. PPR addresses this limitation by diffusing initial suspicion scores across the network topology, allowing latent patterns, collusive groups, and hidden intermediaries to surface naturally.
Formally, the PPR vector $\boldsymbol{r} \in \mathbb{R}^{|V|}$ represents a steady-state distribution of a random walk over the network. At each step, the walk either follows an outgoing transaction with probability $(1-\alpha_{PPR})$ or teleports back to a set of high-suspicion seed nodes with probability $\alpha_{PPR}$:
\begin{equation}
    \boldsymbol{r} = (1 - \alpha_{PPR})P \boldsymbol{r} + \alpha_{PPR} \boldsymbol{v},
\label{eq:ppr}
\end{equation}
where $\alpha_{PPR} \in (0,1)$ is a factor that balances local versus global influence, and $\boldsymbol{v}$ is the personalization vector that concentrates probability mass on a set of known high-risk accounts $Q \subset \mathcal{V}_i$. Specifically:
\begin{equation}
v_i = 
\begin{cases}
    \frac{\hat{\mathbf{y}}_i}{\sum_{j \in Q} \hat{\mathbf{y}}_j}, & \text{if } i \in Q, \\
    0, & \text{otherwise,}
\end{cases}
\end{equation}
where $P_i$ is the initial suspiciousness score of node $i$, derived from prior models, expert labeling, or regulatory flags. 


\subsection{Cross-bank Personalized PageRank}
While PPR effectively amplifies and contextualizes suspiciousness signals within a single financial institution, it faces inherent limitations in the cross-border setting. Because of strict privacy and data-sharing regulations, transaction networks are fragmented across banks and jurisdictions, preventing the propagation of suspicion signals across institutional boundaries. Local biases in financial markets, documented in contexts such as crowdfunding~\parencite{ni2025xinyang}, suggest that institutions may systematically underweight cross-border signals, further motivating the need for mechanisms like Cross-bank PPR that explicitly integrate geographically distributed intelligence. As a result, money-laundering activities spanning multiple countries may appear incomplete or innocuous when analyzed in isolation as each country observes only a fragment of the global laundering pattern.  

The effectiveness of PPR-based detection depends fundamentally on the separability of money-laundering groups within the network topology—a property we formalize in Theorem \ref{thm:detectability}. To enhance detectability in fragmented settings, we introduce Cross-bank PPR, which incorporates cross-institutional intelligence while maintaining privacy constraints. Specifically, we modify the standard PPR formulation to incorporate the cross-bank suspicious signal from other countries into PPR design. 
Formally, by incorporating the cross-bank component $\boldsymbol{v}_{cross}$, the formulation of PPR can be updated as follows:
\begin{align}
    \boldsymbol{r} &= (1 - \alpha_{PPR})P \boldsymbol{r} + \alpha_{PPR} (\boldsymbol{v}_{local} + \label{PPR_weight} \boldsymbol{v}_{cross}), \\
    P &= D^{-1}(A + \hat{A})  
\label{eq:ppr}
\end{align}
where $\hat{A}_{ij}=\hat{y}_{ij}$ is the prediction of the edge $\mathcal{E}_{ij}$ and $\hat{A}_{ij}=0$ if the edge $\mathcal{E}_{ik}$ does not exist, $D$ is the degree matrix of $(A + \hat{A})$,  $\boldsymbol{v}_{local}$ encodes the bank’s internal suspicion sources, and $\boldsymbol{v}_{cross}$ represents external suspicious signals towards the cross-bank transaction shared across banks. Different from traditional PPR, we define $P$ the normalized weighted adjacency matrix that captures transaction intensity and encodes the suspicious score of the edge through $A$ and $\hat{A}$ respectively.
The teleportation factor $\alpha_{PPR}$ thus determines the balance between exploring local transactional neighborhoods and returning to known suspicious nodes. By incorporating the cross-bank component $\boldsymbol{v}_{cross}$, we enable a privacy-preserving exchange of aggregated suspicion information, allowing each bank to indirectly benefit from the intelligence gathered by others without revealing raw transactions. This cross-institutional diffusion enriches each bank’s PPR computation, helping surface hidden multi-bank laundering rings that would otherwise remain undetected due to regulatory isolation. In this way, the cross-bank PPR formulation combines fragmented data while maintaining compliance, substantially enhancing both the interpretability and completeness of risk detection.

In practice, the Cross-bank PPR vector can be computed iteratively as:
\begin{align}
    \boldsymbol{r}^{(t+1)} &= (1 - \alpha_{PPR})P \boldsymbol{r}^{(t)} + \alpha_{PPR} (\boldsymbol{v}_{local} + \boldsymbol{v}_{cross})  \\
    \boldsymbol{r}_{\text{n}} &= \frac{\boldsymbol{r}^{(t+1)}}{\text{max}(\boldsymbol{r}^{(t+1)})}, 
\end{align}
which converges to a unique solution under mild conditions. Since the largest score in $\boldsymbol{r}^{(t+1)}$ corresponds to the most suspicious seed node with the highest likelihood of being anomalous, we normalize $\boldsymbol{r}^{(t+1)}$ by dividing each node’s score by the maximum value in $\boldsymbol{r}^{(t+1)}$, thereby converting the raw scores into anomaly probabilities. Using a predefined threshold, nodes with sufficiently high probabilities are identified as suspicious. Intuitively, in the PPR, nodes that are both structurally proximate to and frequently transacting with high-suspicion accounts accumulate higher steady-state probabilities. The PPR spreads risk signals through the network: even accounts not directly involved with known suspicious actors may receive elevated scores if they participate in multi-step chains, loops, or intermediary hubs commonly used in layering and integration stages of money laundering. Notice that high-ranking nodes can be traced through their most probable propagation paths back to seed nodes, offering analysts narrative explanations such as “Account $v$ connects to seeded node $Q$ through a series of small but structured transfers.” These interpretable trails facilitate compliance documentation and regulatory auditing. 
Overall, PPR provides a mathematically grounded and operationally transparent mechanism to uncover money laundering rings by diffusing known suspicion through the transactional network. By integrating graph-theoretic reasoning with domain-aware weighting and temporal adaptation, it bridges the gap between static anomaly scoring and dynamic, relationship-centric risk discovery, thus improving both the precision and explainability of anti-money-laundering analytics.
 
To obtain a more comprehensive view of money-laundering activities across multiple countries, we merge clustering results from different national transaction networks into a unified framework. We start by identifying the cluster dictionary with the largest number of accounts as the initial reference. For each cluster in this reference dictionary, we iteratively examine clusters from other country-specific dictionaries to detect overlapping accounts. When overlaps are found, the clusters are merged to form a larger, consolidated cluster, ensuring that duplicated accounts are avoided. This process is repeated iteratively across all clusters, further merging any partially overlapping clusters to capture extended relational structures. By combining clusters in this way, the resulting merged dictionary integrates patterns spanning multiple financial institutions and geographies, allowing us to recover more complete money-laundering behaviors that may be fragmented across different national datasets. The merged clusters are then cross-validated against known laundering attempts to highlight clusters that contain multiple suspicious accounts, enhancing the interpretability and practical utility of the results.

\subsection{Theoretical Insights of Using PPR for Grouping Money Laundering Behaviors}
Having introduced Cross-bank PPR and demonstrated its empirical effectiveness, we now provide theoretical foundations for understanding when and why PPR-based methods can successfully identify money-laundering groups. We show that Personalized PageRank (PPR) initialized from a malicious seed node concentrates its probability mass inside the laundering group whenever the within-group connectivity is sufficiently stronger than the background connectivity, following the analysis by Avrachenkov and Andersen et al.~\parencite{avrachenkov2015pagerank, andersen2006local}. Let $\mathcal{G}=(\mathcal{V},\mathcal{E})$ denote a random graph  generated by a two-block Stochastic Block Model (SBM) with $n$ nodes (we omit the graph index in this theoretical analysis for notation simplicity) and a planted laundering group $S^\star\subset V$ of size $s$. Edges are sampled independently as
\begin{equation}
\Pr[(u,v)\in \mathcal{E}] =
    \begin{cases}
    p_{\mathrm{in}}, & u,v\in S^\star,\\[0.2em]
    p_{\mathrm{out}}, & \text{otherwise},
    \end{cases}
\end{equation}
Let $A$ be the adjacency matrix, $D$ the degree diagonal, and $P=D^{-1}A$ the random-walk transition operator. Given $\alpha\in(0,1)$ and a one-hot seed vector $\mathbf{e}_{s_0}$ (corresponding to a known malicious seed $s_0\in S^\star$), the Personalized PageRank vector is
\begin{equation}
    \mathbf{r} =(1-\alpha)\,(I-\alpha P)^{-1}\mathbf{e}_{s_0}.
\end{equation}
Denote by $\bar P$ the mean-field (two-block averaged) transition matrix from the expected graph, whose block-averaged version is
\begin{equation}
    M =
    \begin{pmatrix}
        p_{\mathrm{in}} & p_{\mathrm{out}}\\
        p_{\mathrm{out}} & p_{\mathrm{out}}
    \end{pmatrix}.
\end{equation}
We write $\mu_{\mathrm{in}}$ and $\mu_{\mathrm{out}}$ for the mean-field average PPR per node inside and outside the planted group:
\begin{equation}
    \label{mu_in_and_out}
    \mu_{\mathrm{in}} =\frac{(1-\alpha)}{s}\big[(I-\alpha M)^{-1}\big]_{11}, \qquad
    \mu_{\mathrm{out}} =\frac{(1-\alpha)}{n-s}\big[(I-\alpha M)^{-1}\big]_{21},
\end{equation}
and define the \emph{mean-field gap}: $\Delta_{\mathrm{mean}}=\mu_{\mathrm{in}}-\mu_{\mathrm{out}}$.





\begin{theorem}[Detectability of PPR under Planted Laundering Group]
\label{thm:detectability}
Let $G\sim\mathrm{SBM}$ $(n,s,p_{\mathrm{in}},p_{\mathrm{out}})$ with $p_{\mathrm{in}}>p_{\mathrm{out}}$ and assume that degrees concentrate around their expectations (i.e., Lemma~\ref{lem:concentration} holds). Let $\mathbf{r}$ be the Personalized PageRank vector seeded at a known malicious account $s_0\in S^\star$ with a constant $\alpha\in(0,1)$, and $\varepsilon_{\infty}$ denote the entrywise perturbation scale arising from stochastic fluctuations of $P$ around the mean-field transition matrix $\bar P$. There exist constants $C,c>0$ such that, with probability at least $1-n^{-c}$, if the mean-field gap satisfies
\[
\Delta_{\mathrm{mean}} \;\ge\; C\cdot\varepsilon_{\infty},
\]
then 
\begin{enumerate}
    \item the average PPR score on $S^\star$ exceeds that on $V\setminus S^\star$ by at least $\tfrac{1}{2}\Delta_{\mathrm{mean}}$;
    \item ordering nodes by normalized PPR $r(v)/d(v)$ and selecting the prefix with smallest conductance recovers a subset $\widehat S$ with $|\widehat S\cap S^\star|\ge\gamma s$ for some constant $\gamma\in(0,1)$.
\end{enumerate}
In the standard regime $p_{\mathrm{out}}\gtrsim \log n/n$, the detectability condition simplifies to
\[
s\,(p_{\mathrm{in}}-p_{\mathrm{out}}) \;\gtrsim\; C'\sqrt{n p_{\max}\log n},
\]
where $p_{\max}=\max\{p_{\mathrm{in}},p_{\mathrm{out}}\}$. 
\end{theorem}

\vspace{0.5em}
\noindent\textit{Proof.}
See Appendix~\ref{app:proofs}.

Theorem~\ref{thm:detectability} establishes a detectability condition for identifying a money-laundering group using Personalized PageRank (PPR). Intuitively, PPR captures money-laundering behavior because it \emph{propagates the suspiciousness signal} from a known malicious seed node through the transactional network. Nodes densely connected to the seed, likely participating in group laundering, accumulate higher PPR scores, while weakly connected nodes receive lower scores. The theorem formalizes this by showing that a well-separated group will stand out in the PPR ranking.
In our proposed cross-bank PPR, we modify the transition matrix in PPR by assigning higher weights to edges more likely associated with laundering transactions as defined in Equation~\ref{PPR_weight}, which further amplifies scores for nodes strongly connected via suspicious interactions. Specifically, weighted edges increase the transition probability along suspicious links, causing PPR to concentrate more on nodes involved in laundering. This effectively increases \(\mu_{\mathrm{in}}\) relative to \(\mu_{\mathrm{out}}\), enlarging the gap \(\Delta_{\mathrm{mean}}\). Consequently, the detectability threshold is easier to satisfy, and PPR can recover a larger fraction of the laundering group or rank them higher in the graph cut procedure. On the other hand, to ensure that PPR starts with the malicious seed node, we reduce the probability of starting with a normal seed nodes by incorporating the cross-bank suspicious signal from other countries.


\subsection{Label Refinement via Label Propagation}
In the previous subsections, we demonstrated how cross-bank PPR captures group-wise money laundering patterns and provided theoretical insights into the detectability of PPR under planted laundering groups. In Section~\ref{PPR_interpretation}, we further visualize the identified laundering groups to highlight the interpretability of the proposed method. Beyond interpretation, a natural follow-up question arises: \emph{Can the discovered group-wise laundering patterns be further exploited to improve detection performance?} To address this question, we propose a label refinement strategy based on label propagation. This design is motivated by the observation that money laundering activities are inherently group-based and densely connected. Similar to the diffusion mechanism underlying PPR, we hypothesize that malicious signals identified within high-confidence laundering groups can be propagated to neighboring nodes. This is particularly useful for suspicious nodes that receive low confidence scores from the base detection module and may otherwise be overlooked. By propagating malicious signals from confidently identified nodes to their neighbors, we increase the likelihood of recovering such borderline cases.

Formally, given a weighted graph represented by a sparse adjacency matrix, we first construct a row-normalized transition matrix $\mathbf{S}$. Starting from the edge-level prediction $\hat{\mathbf{y}}_{uv}$ produced by the detection module, we compute a node-level malicious score for each account $u$ as
\begin{align}
    \hat{\mathbf{y}}_u = \frac{1}{|\{v \mid (v,u)\in E\}|} \sum_{v:(v,u)\in E} \hat{\mathbf{y}}_{vu},
\end{align}
which corresponds to the average malicious score over all edges incident to node $u$. These node-level scores form the initial label vector $\hat{\mathbf{y}}_{\text{nodes}}$.
Following the same idea of PPR, the label propagation process then iteratively updates the soft node labels according to
\begin{align}
    \mathbf{\hat{R}}^{(t+1)} &= \alpha \mathbf{S}\mathbf{\hat{R}}^{(t)} + (1-\alpha)\hat{\mathbf{y}}_{\text{nodes}}, \\
    \nonumber \mathbf{S} &= \mathbf{\hat{D}}^{-1}(A + A_{\text{PPR}}),
\end{align}
where $\mathbf{\hat{D}}$ is the degree matrix, $A$ denotes the original graph adjacency matrix, and $A_{\text{PPR}}$ encodes the group-wise laundering structures identified by cross-bank PPR. The parameter $\alpha \in (0,1)$ controls the strength of propagation. Iterations continue until convergence, measured by the $\ell_1$-norm difference between successive estimates, or until a predefined maximum number of iterations is reached. The resulting node labels are normalized to ensure numerical stability and rescale the malicious score within the range of [0,1]:
\begin{align}
    \mathbf{\hat{R}}^{(t+1)} \leftarrow \frac{\mathbf{\hat{R}}^{(t+1)}}{\max(\mathbf{\hat{R}}^{(t+1)})}.
\end{align}

Since $\mathbf{\hat{R}}$ represents node-level malicious scores, we need to convert them back to edge-level signals by assigning each edge $(u,v)$ the maximum score of its incident nodes,
\begin{align}
    \mathbf{\bar{R}}_{uv} = \max\left(\mathbf{\hat{R}}^{(t+1)}_u, \mathbf{\hat{R}}^{(t+1)}_v\right).
\end{align}
Finally, the original edge prediction $\hat{\mathbf{y}}_{uv}$ is refined by incorporating the propagated malicious signal:
\begin{align}
    \hat{\mathbf{y}}_{uv} \leftarrow (1-\alpha_{\text{lp}})\hat{\mathbf{y}}_{uv} + \alpha_{\text{lp}}\mathbf{\bar{R}}_{uv},
\end{align}
where $\alpha_{\text{lp}}$ balances the contribution between the base detector and the propagated group-level information.


\section{Hierarchical Adaptive Decision-Making for Financial Security}
\subsection{Hierarchical Adaptive Decision-Making module}
Regulatory authorities face clear-cut cases where intervention is straightforward: transactions flagged with very high confidence as suspicious should be immediately frozen. However, a critical challenge arises with transactions that receive low or moderate suspicion scores. Simply ignoring them risks missing coordinated laundering activities that unfold across multiple accounts and institutions, while indiscriminately freezing low-confidence transactions can lead to excessive disruption and unnecessary false positives. To address this challenge, we propose a \emph{Hierarchical Adaptive Decision-Making module} that integrates local threshold-based reinforcement learning with a global coordinator to ensure the consistency between local and global policies. The module treats the outputs of the graph-based detection as initial probabilistic labels for each transaction, providing prior knowledge of suspicious activity. Then, it captures both short-term (local bank-level) decision feedback and long-term (global system-level) optimization of cooperative anti-laundering policies.

Each bank $b$ manages a local transaction network $\mathcal{G}^b=(\mathcal{V}^b,\mathcal{E}^b)$, where each transaction $\mathcal{E}_{uv}^b$ carries an anomaly probability $\hat{\mathbf{y}}_{uv}$.  
At time $t$, the bank maintains an intervention threshold $\tau_b(t)$ and decides whether to intervene:
\begin{equation}
    a_{uv}^b:= \text{Intervene on } \mathcal{E}_{uv}^b \iff \hat{\mathbf{y}}_{uv} \geq \tau_b(t).
\end{equation}

To adaptively determine the threshold, we formulate the problem of adjusting thresholds as a reinforcement learning (RL) problem~\parencite{li2017deep, szepesvari2022algorithms, ladosz2022exploration}. At time $t$, the agent observes state $s_{uv}^b$, selects an action $a_{uv}^b$ corresponding to a threshold $\tau_b(t)$, and receives a reward $r_t^b$ that measures the quality of decisions across all edges $\mathcal{E}^{b}_t$ at the time $t$:
\begin{align}
    r_t^b &= \sum_{(u,v)\in \mathcal{E}^{b}_t} R(s_{uv}^b, a_{uv}^b, C_{uv}, y_{uv}) \\
    R(s_{uv}^b, a_{uv}^b, C_{uv}, y_{uv}) &= 
    \begin{cases}
        \alpha_{r1}\log(C_{uv}), & \text{if } y_{uv}=1 \text{ and } a_{uv}^b = \text{Freeze}   \\
        \alpha_{r2}\log(C_{uv}), & \text{if } y_{uv}=1 \text{ and } a_{uv}^b = \text{Monitor} \\
        -\alpha_{r3}\log(C_{uv}), & \text{if } y_{uv}=1 \text{ and } a_{uv}^b = \text{No Intervention} \\
        -\alpha_{r4}, & \text{if } y_{uv}=0 \text{ and } a_{uv}^b = \text{Freeze}   \\
        -\alpha_{r5}, & \text{if } y_{uv}=0 \text{ and } a_{uv}^b = \text{Monitor} \\
        \alpha_{r6}, & \text{if } y_{uv}=0 \text{ and } a_{uv}^b = \text{No Intervention} \\
     \end{cases}
\end{align}
% where $\text{IT}=\{\text{delay}, \text{monitor}\}$.
where $\text{C}_{uv}$ is the cost of missing effective actions that should be deployed (which could be the amount of financial loss in this transaction), $y_{uv}$ is the label showing whether this transaction is illicit, $\alpha_{r1}$, $\alpha_{r2}$, $\alpha_{r3}$, $\alpha_{r4}$, $\alpha_{r5}$ and $\alpha_{r6}$ are positive hyperparameters balancing the importance of different actions in six scenarios. Specifically, when a transaction is malicious (i.e.,  $(y_{uv}=1$), we reward the model if the model begins to intervene in this transaction either by freezing or monitoring this transaction, with $\alpha_{r1} > \alpha_{r2} > \alpha_{r3}$ showing the different weight for three actions. Otherwise, we reduce its rewards if no intervention is deployed. When a transaction is normal activity, we reward the model if it does not intervene in this normal behavior. Otherwise, we penalize the model at different degrees with $\alpha_{r4} \geq \alpha_{r6} >\alpha_{r5}$. This reward design encourages high-probability suspicious transactions to trigger appropriate interventions, penalizes unnecessary actions on benign transactions, and incorporates operational cost through both the constant penalties and the cost-aware logarithmic terms. The goal of the local policy $\pi_\theta^b(a|s_{uv}^b)$ is to maximize the expected cumulative reward:
\begin{equation}
    J(\pi_\theta^b) = \mathbb{E}_{\pi^b}\!\Big[\sum_{t=0}^{\infty} \gamma^t r_t^b \Big].
\end{equation}

Our dynamic decision-making framework draws on the tradition of modeling sequential choice under uncertainty, where agents make decisions that affect both immediate outcomes and future states~\parencite{mehta2017ni}. To address the negative impact of bank isolation issue, we propose a global-local coordinator to ensure the consistency between local and global policies. Specifically, local institutions adaptively respond to evolving transaction patterns, while the global coordinator aligns their behavior toward collective objectives. The global policy $\pi_\phi^{\text{global}}$ updates local thresholds using soft coordination:
\begin{equation}
    \tau_b(t+1) \leftarrow \tau_b(t) + \eta_g \big(\bar{\tau}(t) - \tau_b(t)\big), \quad 
    \bar{\tau}(t) = \sum_{b=1}^{B} w_b \tau_b(t), 
\end{equation}
where $\omega_b$ represents the relative institutional importance or transaction volume weight. To ensure the consistency between local and global policies, we define the coupling constraint as:
\begin{align}
    \mathcal{L}_{\text{couple}} &= 
    \sum_{b=1}^B \big\| \tau_b(t) - \bar{\tau}(t) \big\|^2
    + \xi \sum_{b=1}^B D_{\mathrm{KL}}\!\big(\pi_\theta^b \,||\, \pi_\phi^{\text{global}}\big), 
\end{align}
where the first term enforces consistency between local and global thresholds, and the second term aligns local policies with global intent through a KL-divergence regularization weighted by $\xi$. 
The joint optimization objective becomes:
\begin{equation}
    \max_{\{\pi_\theta^b\},\, \pi_\phi^{\text{global}}}
    \sum_{b=1}^B J(\pi_\theta^b)
    - \lambda_{\text{c}} \mathcal{L}_{\text{couple}},
\end{equation}
where $\lambda_{\text{c}}$ controls the strength of policy coupling.
The proposed hierarchical adaptive decision-making framework establishes a principled coordination structure for multi-institution financial security systems. Local RL agents autonomously learn transaction-level thresholds for each bank, while a global coordinator ensures cooperative, privacy-preserving, and regulation-compliant adaptation.  

